\documentclass[11pt]{article}
\usepackage[a4paper,margin=25mm]{geometry}
\usepackage{amsmath,amssymb,mathtools}
\usepackage{fontspec,unicode-math}
\setmainfont{Latin Modern Roman}
\setsansfont{Latin Modern Sans}
\setmonofont{Latin Modern Mono}
\setmathfont{Latin Modern Math}
\usepackage{xurl,hyperref}
\hypersetup{colorlinks=true,urlcolor=blue,linkcolor=blue}
\newcommand{\sref}[2]{\hyperlink{source:#1:#2}{[S#2]}}
\newcommand{\cataloguescope}[1]{\paragraph{Mathematical question.}#1}
\begin{document}
\section{Recognizing genus-three Jacobians over the base field}\label{statement}

\subsection{Definitions and mathematical statement}

\paragraph{Shared notation.}$\mathbb N=\{1,2,\ldots\}$, and $\mathbb Z,\mathbb Q,\mathbb R,\mathbb C$ denote the integers, rationals, reals and complex numbers. Any other conventions are specified locally.


For an abelian variety $A/\mathbb F_q$, its Weil polynomial is $f_A(T)=\det(T-\operatorname{Frob}_q\mid V_\ell A)$, where $\ell$ is a prime different from the characteristic, $V_\ell A=(\varprojlim_r A[\ell^r](\overline{\mathbb F}_q))\otimes_{\mathbb Z_\ell}\mathbb Q_\ell$, and $\operatorname{Frob}_q$ is the $q$-power Frobenius. It has degree twice $\dim A$. An isogeny is a surjective homomorphism of abelian varieties with finite kernel. The Jacobian $\operatorname{Jac}(C)$ of a smooth projective geometrically connected curve $C$ is its abelian variety of degree-zero divisor classes. A principal polarization $\lambda:A\to A^\vee$ is the isomorphism to the dual abelian variety induced by an ample line bundle. A Jacobian is regarded with its canonical principal polarization $\lambda_C$.
\paragraph{Statement recorded in the supplied source.}
For a field $k$ and a principally polarized abelian threefold $(A,\lambda)/k$, determine a criterion for the existence of a smooth projective geometrically connected genus-three curve $C/k$ and an isomorphism of polarized abelian varieties
\[
(A,\lambda)\simeq(\operatorname{Jac}(C),\lambda_C)
\]
defined over $k$. Address the base-field descent obstruction as well as the geometric principal-polarization condition. \sref{statement}{2}

\subsection{Short English statement}
Recognize when a principally polarized abelian threefold is the Jacobian of a curve over its own base field.
\subsection{Sources}
\begin{description}
\item[\hypertarget{source:statement:1}{S1}] ULAM.AI and the UnsolvedMath Contributors, UnsolvedMath: A Curated Collection of Open Mathematics Problems, catalogue entry. CC BY 4.0. \url{https://huggingface.co/datasets/ulamai/UnsolvedMath}.
\item[\hypertarget{source:statement:2}{S2}] Everett W. Howe, \emph{Low-genus curves over finite fields: Problems and variations}, IPAM Number Theory and Cryptography: Open Problems workshop (9--13 October 2006), PDF pages 31--34, ``A half-open problem. \url{https://helper.ipam.ucla.edu/publications/SCWS1/SCWS1_6237.pdf}.
\end{description}
\label{end:statement}
\end{document}
